\documentclass{ximera}

\title{Introductie}
\author{Mila Vervoort}
\license{CC: 0}         % replace with an appropriate license, or set it in xmPreamble

\begin{document}
\begin{abstract}
    Samenvatting van de theorie over reactiekinetiek.
\end{abstract}
\maketitle
\label{xim:Introductie analysemethode}

\section*{Analyse en verwerking van kinetische metingen}

Om een reactor van een bepaald type te ontwerpen, is een uitdrukking nodig voor de chemische reactiesnelheid. In combinatie met de massabalans en de stoichiometrische relaties laat deze snelheidsvergelijking toe om een isotherm reactiesysteem te dimensioneren.

Eigenschappen van chemische reacties zijn vaak ongekend en niet eenvoudig te verkrijgen. Soms, voornamelijk voor elementaire gasfase reacties, zijn er voorspellende methoden mogelijk die nauwkeurige schattingen kunnen geven. De meest gebruikte methode om chemische kinetica te bepalen is echter via
experimentele metingen. In dit hoofdstuk behandelen we de analyse
en verwerking van experimentele metingen om kinetische parameters te bepalen.

We gaan uit van een functionele vorm voor de reactiesnelheid:
\[-R_A = f(C_A)g(T)\]

Voor kinetische metingen worden doorgaans twee types reactoren gebruikt:
\begin{itemize}
\item \textbf{Batchreactoren}, voornamelijk voor de studie van homogene reacties;
\item \textbf{Continue reactoren} (kuipreactoren of buisreactoren), die vaak gebruikt worden wanneer de interpretatie van meetgegevens eenvoudiger is onder stationaire condities.
\end{itemize}


\subsection*{Kinetische metingen in een batch reactor}
\label{Kinetische metingen in een batch reactor}

Deze methode wordt vooral gebruikt voor homogene reacties. We veronderstellen een isotherme reactor met constant volume.

Om een uitdrukking voor de reactiesnelheid te bepalen, worden concentratiemetingen op verschillende tijdstippen uitgevoerd. Soms wordt de concentratie gekoppeld aan andere meetbare eigenschappen, zoals: thermische conductiviteit, viscositeit, spectrometrie of de totaaldruk.

\begin{example} \textbf{Concentratiebepaling via drukmetingen}

In dit voorbeeld wordt aangetoond hoe de concentraties in een batchreactor voor gasfasereacties eenvoudig bepaald kunnen worden door de totale druk en de partieel druk te meten.

Beschouw een reactie met stoichiometrie:
\[
aA + bB \rightarrow rR + sS
\]
De molaire hoeveelheden van de componenten gedurende de reactie zijn:
\[
\begin{array}{c|ccccc}
 & A & B & R & S & I \\ \hline
t=0 & N_{A0} & N_{B0} & N_{R0} & N_{S0} & N_{I0} \\
t=t & N_{A0}-ax & N_{B0}-bx & N_{R0}+rx & N_{S0}+sx & N_{I0}
\end{array}
\]
waar $x$ de reactiegraad is.\\
De totale hoeveelheid mol op tijdstip $t$ wordt dan:
\[
N = N_0 + x \Delta n \text{   met   } \Delta n = (r + s - a - b)
\]
De concentratie van component $A$ is
\[
C_A = \frac{N_A}{V} = \frac{N_{A0} - ax}{V}.
\]
Voor een ideaal gas geldt de ideale gaswet: 
\[V p_A = N_A R T\]
Omgevormd naar concentratie wordt dit:
\[C_A = \frac{N_A}{V} = \frac{p_A}{RT}\]
De concentratie in functie van de totale druk en de systeemdruk wordt dan:
\[C_A = \frac{p_A}{RT} \] met 
\[p_A = p_{A0} - \frac{a}{\Delta n}(p-p_0)\]
\end{example}

Wanneer kinetische metingen de concentratie $C_A$ op verschillende tijdstippen $t$ opleveren, kan de reactiekinetiek hieruit worden afgeleid. 

Hiervoor kunnen verschillende methoden worden gebruikt:
\begin{itemize}
    \item Integrale analysemethode
    \item Differentiële analysemethode
    \item Methode der beginsnelheden
    \item Methode der halfwaardetijden
    \item Kleinste kwadraten benadering
\end{itemize}

\subsection*{Kinetische metingen in een continue reactor}

\end{document}